\documentclass[11pt,a4paper]{article}
\usepackage[margin=19mm]{geometry}
\usepackage[T1]{fontenc}
\usepackage{lmodern,amsmath,amssymb,mathtools,booktabs,array,xcolor,hyperref}
\hypersetup{hypertexnames=false,colorlinks=true,linkcolor=teal,urlcolor=teal,pdftitle={LQG equation atlas: formula review},pdfauthor={Prepared for Antoine Dierckx}}
\definecolor{ink}{HTML}{173D4A}
\newcommand{\ket}[1]{\lvert #1\rangle}
\newcommand{\bra}[1]{\langle #1\rvert}
\newcommand{\Inv}{\operatorname{Inv}}
\newcommand{\Contr}{\operatorname{Contr}}
\newcommand{\Tr}{\operatorname{Tr}}
\newcommand{\Vol}{\operatorname{vol}}
\newcommand{\eq}[2]{\begin{equation}#2\tag{#1}\end{equation}}
\setlength{\parindent}{0pt}
\setlength{\parskip}{5pt}
\setlength{\emergencystretch}{2em}
\allowdisplaybreaks[2]
\begin{document}
{\color{ink}\Huge LQG equation atlas}\par
{\Large Proposed formula catalogue}\par
\textbf{Review draft for Antoine Dierckx \quad | \quad 6 September 2026}

This document restores the equation catalogue from the chat answer. Its labels are stable feedback targets: for example, ``change (G5)'' or ``expand (D6) further.'' It is a proposal for scientific review, not a claim that every convention-dependent translation has already passed implementation tests. No website code or deployment is included.

\textbf{Proposed common convention.} Use Speziale's Lorentz-matrix phase convention, which makes the usual booster functions real. Express Toller matrices in the same basis. The conversion is trivial on the minimal blocks used directly in the EPRL and causal vertices; the general off-diagonal Feynman formula needs an explicit phase ratio.

\textbf{Corrections to the original brief.} The normalized one-loop state is $\chi_j$, not $\sqrt{d_j}\chi_j$. Booster dimension factors depend on intertwiner normalization. Toller matrices do not obey the Wigner composition law, so the ordinary booster decomposition cannot be transferred by the replacement $d\mapsto t$.

\textbf{Still to fix before implementation.} The exact named $15j$ mapping and its orientation phases must be checked against the supplied TikZ port order; the fully indexed tensor contraction is the reference. The AL/RS volume comparison retains explicit regularization constants. The first-order action's orientation and the connection/holonomy endpoint convention must be locked together. These points are marked in the relevant sections.

\begingroup
\footnotesize
\setlength{\parskip}{0pt}
\tableofcontents
\endgroup
\newpage
\section{Conventions and the three initial equations}
Keep $\hbar$ explicit, set $c=1$, and initially take $\gamma>0$:
\eq{C1}{\eta_{IJ}=\operatorname{diag}(-1,1,1,1),\quad\kappa=8\pi G,\quad\ell_P^2=G\hbar,\quad\ell_0^2=\kappa\gamma\hbar.}
\eq{C2}{J_i=\frac{\sigma_i}{2},\quad\tau_i=-iJ_i,\quad[J_i,J_j]=i\epsilon_{ij}{}^kJ_k,\quad d_j=2j+1.}
\eq{C3}{D^j_{mn}(h)=\bra{j,m}D^j(h)\ket{j,n},\qquad\int_{SU(2)}dh=1.}
Clebsch--Gordan coefficients use Condon--Shortley phases. The duality tensor is
\eq{C4}{\epsilon^{(j)}_{mn}=(-1)^{j-m}\delta_{m,-n},\qquad\epsilon^{(j)T}=(-1)^{2j}\epsilon^{(j)}.}
These conventions determine signs from wire reversals and leg permutations.

The initial screen shows:
\eq{T1}{\boxed{\mathcal H_\Gamma=L^2[SU(2)^L/SU(2)^N].}}
\eq{T2}{\boxed{[\hat X_l^i,\hat X_{l'}^j]=i\ell_0^2\delta_{ll'}\epsilon^{ij}{}_k\hat X_l^k,\qquad\hat{\mathbf X}_l=\ell_0^2\hat{\mathbf L}_l.}}
\eq{T3}{\boxed{A_v(\psi_v)=\left(P_{SL(2,\mathbb C)}Y_\gamma\psi_v\right)(\mathbb 1).}}
The dimensionless algebra and the area and volume formulas are immediate children of (T2). In (T3), $P$ denotes formal noncompact group averaging, evaluated with the redundant vertex integration removed. It is not an ordinary bounded orthogonal projector on an $L^2$ space.

\section{Classical metric and first-order branches}
\eq{A1}{g_{\mu\nu}=\eta_{IJ}e^I_\mu e^J_\nu.}
\eq{A2}{S_{\rm EH}[g]=\frac1{2\kappa}\int_Md^4x\sqrt{-g}(R-2\Lambda).}
\eq{A3}{(\star B)_{IJ}=\frac12\epsilon_{IJ}{}^{KL}B_{KL},\qquad F^{IJ}=d\omega^{IJ}+\omega^I{}_K\wedge\omega^{KJ}.}
Choose the tetrad orientation consistently with (A2). The proposed first-order convention is
\eq{A4}{S[e,\omega]=\frac1{2\kappa}\int_M\left[\star(e\wedge e)_{IJ}+\gamma^{-1}(e\wedge e)_{IJ}\right]\wedge F^{IJ}-\frac{\Lambda}{\kappa}\int_M\Vol_e,}
\eq{A5}{\Vol_e=\frac1{4!}\epsilon_{IJKL}e^I\wedge e^J\wedge e^K\wedge e^L.}
The torsion-free reduction is
\eq{A6}{T^I=D_\omega e^I,\qquad T^I=0\quad\text{(vacuum, nondegenerate tetrad, real $\gamma$).}}
Removing the Holst term does not automatically preserve every subsequent $\gamma$-dependent quantization formula.

Optional action densities are
\eq{A7}{\begin{aligned}
\mathcal N_{\rm NY}&=d(e_I\wedge T^I)=T_I\wedge T^I-e_I\wedge e_J\wedge F^{IJ},\\
\mathcal P&=F^{IJ}\wedge F_{IJ},\\
\mathcal E&=\tfrac12\epsilon_{IJKL}F^{IJ}\wedge F^{KL}.
\end{aligned}}
Their coefficients are independent parameters. A smooth non-null metric-boundary example is
\eq{A8}{S_{\rm GHY}=\frac1\kappa\int_{\partial M}d^3x\,\varepsilon_n\sqrt{|h|}\,K,\qquad\varepsilon_n=n_\mu n^\mu=\pm1.}
This term must be labelled with its boundary conditions; null boundaries and joints require separate formulas. The undeformed EPRL model below has $\Lambda=0$. A classical cosmological-constant switch must not silently change that model.

\section{Canonical variables and constraints}
\eq{K1}{ds^2=-N^2dt^2+q_{ab}(dx^a+N^adt)(dx^b+N^bdt).}
In time gauge:
\eq{K2}{E^a_i=\tfrac12\epsilon^{abc}\epsilon_{ijk}e_b^je_c^k,\quad E^a_iE^{bi}=q\,q^{ab},\quad q=\det q_{ab}.}
\eq{K3}{A_a^i=\Gamma_a^i(E)+\gamma K_a^i,\qquad K_a^i=K_{ab}e^{bi}.}
\eq{K4}{de^i+\epsilon^i{}_{jk}\Gamma^j\wedge e^k=0.}
\eq{K5}{\{A_a^i(x),E^b_j(y)\}=\kappa\gamma\delta_a^b\delta^i_j\delta^{(3)}(x-y),\qquad\{A,A\}=\{E,E\}=0.}
\eq{K6}{G_i=\partial_aE^a_i+\epsilon_{ij}{}^kA_a^jE^a_k\approx0.}
\eq{K7}{D[\mathbf N]=\frac1{\kappa\gamma}\int_\Sigma d^3x\,N^a(F_{ab}^iE^b_i-A_a^iG_i)\approx0.}
\eq{K8}{\begin{aligned}
H[N]={}&\frac1{2\kappa}\int_\Sigma d^3x\,N\frac{E^a_iE^b_j}{\sqrt q}
\left[\epsilon^{ij}{}_kF_{ab}^k-2(1+\gamma^2)K_{[a}^iK_{b]}^j\right]\\
&+\frac\Lambda\kappa\int_\Sigma d^3x\,N\sqrt q\approx0.
\end{aligned}}
Antisymmetrization includes $1/2$. This is a classical constraint; the atlas must not imply that its quantum regularization is unique.

\section{Holonomies, fluxes and endpoint algebra}
\eq{H1}{h_e[A]=\mathcal P\exp\left(\int_e A_a^i\tau_i\,dx^a\right),\qquad E_i(S)=\frac12\int_S\epsilon_{abc}E_i^a\,dx^b\wedge dx^c.}
The ordinary continuum flux $E_i(S)$ and a gauge-covariant graph flux $\mathbf X_l$ are distinct constructions. The path-ordering convention must agree with the source/target matrix convention used below.
\eq{H2}{\begin{aligned}
(\hat L_s^if)(h)&=-i\left.\frac d{dt}f(e^{t\tau_i}h)\right|_{t=0},\\
(\hat L_t^if)(h)&=-i\left.\frac d{dt}f(he^{-t\tau_i})\right|_{t=0}.
\end{aligned}}
\eq{H3}{[\hat L_s^i,\hat L_s^j]=i\epsilon^{ij}{}_k\hat L_s^k,\quad[\hat L_t^i,\hat L_t^j]=i\epsilon^{ij}{}_k\hat L_t^k,\quad[\hat L_s^i,\hat L_t^j]=0.}
\eq{H4}{[\hat L_s^i,\hat h]=-J_i\hat h,\qquad[\hat L_t^i,\hat h]=\hat hJ_i.}
Both endpoint generators are treated as outgoing at their respective nodes:
\eq{H5}{\hat G_n^i\Psi=\left(\sum_{l:s(l)=n}\hat L_{s,l}^i+\sum_{l:t(l)=n}\hat L_{t,l}^i\right)\Psi=0.}
\eq{H6}{\hat{\mathbf L}_t=-\operatorname{Ad}_{h^{-1}}\hat{\mathbf L}_s.}
Diagrams should show this endpoint convention whenever these equations appear.

\section{Hilbert space and Peter--Weyl expansion}
\eq{S1}{\widetilde{\mathcal H}_\Gamma=L^2(SU(2)^L,dh^L),\qquad\mathcal H_\Gamma=\widetilde{\mathcal H}_\Gamma^{\,SU(2)^N}.}
\eq{S2}{L^2(SU(2))\simeq\widehat{\bigoplus}_{j\in\frac12\mathbb N_0}V_j\otimes V_j^*.}
\eq{S3}{\int dh\,\overline{D^j_{mn}(h)}D^{j'}_{m'n'}(h)=\frac{\delta_{jj'}\delta_{mm'}\delta_{nn'}}{d_j}.}
\eq{S4}{\begin{aligned}
f(h)&=\sum_{jmn}f^j_{mn}\sqrt{d_j}\,D^j_{mn}(h),\\
f^j_{mn}&=\int dh\sqrt{d_j}\,\overline{D^j_{mn}(h)}f(h).
\end{aligned}}
\eq{S5}{\widetilde{\mathcal H}_\Gamma\simeq\widehat{\bigoplus}_{\{j_l\}}\bigotimes_l(V_{j_l}\otimes V_{j_l}^*).}
\eq{S6}{(P_\Gamma f)(\{h_l\})=\int\prod_n du_n\,f(\{u_{s(l)}h_lu_{t(l)}^{-1}\}).}
\eq{S7}{\boxed{\mathcal H_\Gamma\simeq\widehat{\bigoplus}_{\{j_l\}}\bigotimes_n\Inv_{SU(2)}\left[\bigotimes_{l:s(l)=n}V_{j_l}\otimes\bigotimes_{l:t(l)=n}V_{j_l}^*\right].}}
\eq{S8}{\Psi_{\Gamma,j_l,i_n}(\{h_l\})=\Contr_\Gamma\left[\bigotimes_l\sqrt{d_{j_l}}D^{j_l}(h_l)\otimes\bigotimes_n i_n\right].}
The node tensors are normalized. $\Contr_\Gamma$ unfolds into magnetic sums and the duality tensors (C4), with appropriate dual node tensors. This is a fixed-graph Gauss-invariant kinematical Hilbert space, not the fully constrained physical space.

\section{Intertwiners and changes of coupling}
\eq{I1}{i_{m_1m_2m_3}=\begin{pmatrix}j_1&j_2&j_3\\m_1&m_2&m_3\end{pmatrix}.}
\eq{I2}{C^{JM}_{j_1m_1\,j_2m_2}=(-1)^{j_1-j_2+M}\sqrt{2J+1}\begin{pmatrix}j_1&j_2&J\\m_1&m_2&-M\end{pmatrix}.}
\eq{I3}{|j_1-j_2|\le j_3\le j_1+j_2,\quad j_1+j_2+j_3\in\mathbb Z,\quad m_1+m_2+m_3=0.}
For four ordered outgoing legs, use the normalized basis
\eq{I4}{(i_k)_{m_1m_2m_3m_4}=\sum_{\mu=-k}^k\frac{(-1)^{k-\mu}}{\sqrt{d_k}}C^{k\mu}_{j_1m_1\,j_2m_2}C^{k,-\mu}_{j_3m_3\,j_4m_4}.}
\eq{I5}{\begin{gathered}
k\in[|j_1-j_2|,j_1+j_2]\cap[|j_3-j_4|,j_3+j_4],\\
j_1+j_2+k\in\mathbb Z,\qquad j_3+j_4+k\in\mathbb Z.
\end{gathered}}
\eq{I6}{\langle i_k|i_{k'}\rangle=\delta_{kk'},\qquad P_{\Inv}=\sum_k\ket{i_k}\bra{i_k}.}
\eq{I7}{\begin{aligned}
\ket{((j_1j_2)k\,j_3)JM}={}&\sum_{k'}(-1)^{j_1+j_2+j_3+J}\sqrt{d_kd_{k'}}\\
&\times\begin{Bmatrix}j_1&j_2&k\\j_3&J&k'\end{Bmatrix}\ket{(j_1(j_2j_3)k')JM}.
\end{aligned}}
For arbitrary diagram pairings the exact change-of-basis matrix is
\eq{I8}{U_{k'k}=\sum_{\mathbf m}\overline{(i^{\rm new}_{k'})_{\mathbf m}}(i^{\rm old}_{k})_{\mathbf m}.}

\section{Area, angle and volume}
\eq{G1}{\hat A_l=\ell_0^2\sqrt{\hat{\mathbf L}_l^2},\qquad A_j=\ell_0^2\sqrt{j(j+1)}.}
\begin{center}\begin{tabular}{c|ccccc}$j$&$0$&$1/2$&$1$&$3/2$&$2$\\\hline
$A_j/\ell_0^2$&$0$&$\sqrt3/2$&$\sqrt2$&$\sqrt{15}/2$&$\sqrt6$\end{tabular}\end{center}
For transverse intersections away from nodes:
\eq{G2}{\hat A(S)=\ell_0^2\sum_{p\in S\cap\Gamma}\sqrt{\hat{\mathbf L}_p^2}.}
For nonzero outward face normals:
\eq{G3}{\widehat{\cos\theta_{ab}}=\frac{\hat{\mathbf L}_a\cdot\hat{\mathbf L}_b}{\sqrt{\hat{\mathbf L}_a^2}\sqrt{\hat{\mathbf L}_b^2}}.}
If $j_a,j_b$ couple to $k$,
\eq{G4}{\cos\theta_{ab}=\frac{k(k+1)-j_a(j_a+1)-j_b(j_b+1)}{2\sqrt{j_a(j_a+1)j_b(j_b+1)}}.}
The interior dihedral angle of a convex polyhedron is the supplement of this outward-normal angle.
\eq{G5}{\hat Q=\epsilon_{ijk}\hat L_1^i\hat L_2^j\hat L_3^k,\qquad\boxed{\hat V_{\rm tet}=\frac{\sqrt2}{3}\ell_0^3\sqrt{|\hat Q|}.}}
\eq{G6}{\begin{aligned}
\hat Q&=i[\hat{\mathbf L}_1\cdot\hat{\mathbf L}_2,\hat{\mathbf L}_2\cdot\hat{\mathbf L}_3]\\
&=\frac i4[(\hat{\mathbf L}_1+\hat{\mathbf L}_2)^2,(\hat{\mathbf L}_2+\hat{\mathbf L}_3)^2].
\end{aligned}}
The advanced comparison would retain explicit regularization normalizations:
\eq{G7}{\begin{aligned}
\hat V_{{\rm AL},n}&=C_{\rm AL}\ell_0^3\sqrt{\left|\sum_{a<b<c}\epsilon(e_a,e_b,e_c)\hat Q_{abc}\right|},\\
\hat V_{{\rm RS},n}&=C_{\rm RS}\ell_0^3\sum_{a<b<c}\sqrt{|\hat Q_{abc}|}.
\end{aligned}}
The $C$'s record regularization conventions, not new physical couplings. Do not equate these operators to (G5) without an embedding and normalization specification.

\section{Explicit low-spin examples}
The normalized gauge-invariant one-loop state is
\eq{L1}{\psi_j(h)=\chi_j(h)=\Tr D^j(h),\qquad\int dh\,\overline{\chi_j(h)}\chi_{j'}(h)=\delta_{jj'}.}
For eigenvalues $e^{\pm i\theta/2}$ of $h$,
\eq{L2}{\chi_j(h)=\frac{\sin[(2j+1)\theta/2]}{\sin(\theta/2)},}
with apparent singularities interpreted by continuity, and
\eq{L3}{\chi_0=1,\qquad\chi_{1/2}=2\cos(\theta/2),\qquad\chi_1=1+2\cos\theta.}
For $(1/2,1/2,1)$ a normalized invariant is
\eq{L4}{\ket{i}=\frac1{\sqrt3}\left[\ket{++,-1}-\frac{\ket{+-,0}+\ket{-+,0}}{\sqrt2}+\ket{--,+1}\right].}
Gauge-invariant trivalent nodes have zero volume.
\eq{L5}{\Inv(V_{1/2}^{\otimes4})=\operatorname{span}\{\ket0,\ket1\}.}
\eq{L6}{\begin{aligned}
\ket0&=\ket{s}_{12}\ket{s}_{34},\qquad\ket{s}=\frac{\ket{+-}-\ket{-+}}{\sqrt2},\\
\ket1&=\frac{\ket{t_+}\ket{t_-}-\ket{t_0}\ket{t_0}+\ket{t_-}\ket{t_+}}{\sqrt3}.
\end{aligned}}
With the phases in (I4), direct Clebsch--Gordan evaluation gives
\eq{L7}{\hat Q_{\frac12^4}=\frac{\sqrt3}{4}\begin{pmatrix}0&-i\\i&0\end{pmatrix}.}
\eq{L8}{\ket{q_\pm}=\frac{\ket0\pm i\ket1}{\sqrt2},\quad q_\pm=\pm\frac{\sqrt3}{4},\quad V_\pm=\frac{3^{1/4}}{3\sqrt2}\ell_0^3.}
The states have opposite oriented triple products, but the same positive volume.
For four spins $1$, in the $k=0,1,2$ basis:
\eq{L9}{\hat Q_{1^4}=\begin{pmatrix}0&-2i/\sqrt3&0\\2i/\sqrt3&0&-i\sqrt{5/3}\\0&i\sqrt{5/3}&0\end{pmatrix}.}
\eq{L10}{\operatorname{spec}\hat Q=\{0,+\sqrt3,-\sqrt3\},\qquad\operatorname{spec}\hat V=\left\{0,\frac{\sqrt2\,3^{1/4}}3\ell_0^3\right\}.}
The nonzero volume is doubly degenerate. A zero-volume eigenvector is
\eq{L11}{\ket{q_0}=\frac{\sqrt5\ket0+2\ket2}{3}.}

\section{Lorentz representations and simplicity}
\eq{Y1}{[L_i,L_j]=i\epsilon_{ij}{}^kL_k,\quad[L_i,K_j]=i\epsilon_{ij}{}^kK_k,\quad[K_i,K_j]=-i\epsilon_{ij}{}^kL_k.}
\eq{Y2}{\mathcal V^{(\rho,k)}=\widehat{\bigoplus}_{j=|k|}^\infty V_j,\qquad\rho\in\mathbb R,\quad k\in\tfrac12\mathbb Z.}
\eq{Y3}{\mathbf K^2-\mathbf L^2=\rho^2-k^2+1,\qquad\mathbf K\cdot\mathbf L=\rho k.}
\eq{Y4}{Y_\gamma:V_j\longrightarrow\mathcal V^{(\gamma j,j)},\qquad Y_\gamma\ket{j,m}=\ket{\gamma j,j;j,m}.}
\eq{Y5}{Y_\gamma^\dagger Y_\gamma=\mathbb1_{V_j},\qquad Y_\gamma Y_\gamma^\dagger=P_{j,{\rm minimal}}.}
For $j>0$,
\eq{Y6}{P_jK_iP_j=\frac{\rho k}{j(j+1)}L_i.}
Thus the default map gives
\eq{Y7}{Y_\gamma^\dagger K_iY_\gamma=\gamma\frac{j}{j+1}L_i.}
The alternative map gives
\eq{Y8}{\rho=\gamma(j+1),\quad k=j\quad\Longrightarrow\quad Y^\dagger(K_i-\gamma L_i)Y=0.}
This is a projected matrix-element statement, not a strong equality on the whole Lorentz representation. See the primer's canonical-basis and Casimir discussion [1].
The BF bridge, on the chosen gravitational sector, is
\eq{Y9}{S_{\rm BF}=\frac1{2\kappa}\int B_{IJ}\wedge F^{IJ},\qquad B_{IJ}=\star(e\wedge e)_{IJ}+\gamma^{-1}(e\wedge e)_{IJ}.}
``Simplicity off'' must lead to the BF representation measure and amplitudes, not just delete a $Y_\gamma$ symbol.

\section{Vertex: group variables and magnetic indices}
\eq{D1}{K_j(g)=Y_\gamma^\dagger D^{(\gamma j,j)}(g)Y_\gamma,\qquad[K_j(g)]_{mn}=D^{(\gamma j,j)}_{jm,jn}(g).}
Formal Lorentz averaging on a connected vertex boundary is
\eq{D2}{(P_{SL(2,\mathbb C)}\Phi)(\{g_l\})=\int\prod_n dG_n\,\Phi(\{G_{s(l)}g_lG_{t(l)}^{-1}\}).}
Evaluation at the identity uses one integration fewer. For the four-simplex,
\eq{D3}{g_1=\mathbb1,\qquad g_{ab}=g_a^{-1}g_b,\qquad1\le a<b\le5.}
A group-variable kernel, fixing the wedge-holonomy convention, is
\eq{D4}{\begin{aligned}
A_v(\{h_{ab}\})={}&\sum_{\{j_{ab}\}}\left(\prod_{a<b}d_{j_{ab}}\right)\int\prod_{a=2}^5dg_a\\
&\times\prod_{a<b}\Tr_{j_{ab}}\left[K_{j_{ab}}(g_{ab})D^{j_{ab}}(h_{ab})\right].
\end{aligned}}
\eq{D5}{\Tr[K_j(g)D^j(h)]=\sum_{m,n=-j}^j D^{(\gamma j,j)}_{jm,jn}(g)D^j_{nm}(h).}
For fixed labels,
\eq{D6}{A_v(j_{ab},i_a)=\int\prod_{a=2}^5dg_a\;\Contr_{K_5}\left[\bigotimes_{a<b}K_{j_{ab}}(g_{ab})\otimes\bigotimes_ai_a\right].}
The proposed all-outgoing indexed form is
\eq{D7}{\begin{aligned}
A_v(j_{ab},i_a)={}&\int\prod_{a=2}^5dg_a\sum_{\{m_{ab}\}}\left[\prod_a(i_a)_{\{m_{ab}\}_{b\ne a}}\right]\\
&\times\prod_{a<b}\mathcal K^{j_{ab}}_{m_{ab}m_{ba}}(g_{ab}),
\end{aligned}}
\eq{D8}{\mathcal K^j_{mn}(g)=\sum_{p=-j}^j[K_j(g)]_{mp}\epsilon^{(j)}_{pn}.}
Leg order is the diagram's port order. The endpoint-to-dual assignment must be checked when matching (D4) to normalized boundary wavefunctions; boundary dimension factors must not be inferred by inspection.
The composition law exposes the auxiliary spins:
\eq{D9}{D^{(\rho,k)}_{jm,jn}(g_a^{-1}g_b)=\sum_{\ell=|k|}^\infty\sum_{p=-\ell}^{\ell}\overline{D^{(\rho,k)}_{\ell p,jm}(g_a)}D^{(\rho,k)}_{\ell p,jn}(g_b).}
For EPRL, $\ell\ge j$.

\section{Full foam amplitude and gluing}
\eq{F1}{W_{\mathcal C}(j_\partial,i_\partial)=\sum_{\{j_f,i_e\}_{\rm internal}}\prod_f A_f(j_f)\prod_e A_e(j_f,i_e)\prod_vA_v(j_f,i_e).}
A convenient normalized-intertwiner convention is
\eq{F2}{A_f(j_f)=d_{j_f},\qquad A_e=1,}
with boundary factors chosen consistently with gluing. Other conventions redistribute factors among faces, edges, vertices and boundary states.
\eq{F3}{W_{\mathcal C}(\psi)=\sum_{j_\partial,i_\partial}\psi_{j_\partial,i_\partial}W_{\mathcal C}(j_\partial,i_\partial).}
The wedge-variable counterpart has the structure
\eq{F4}{W_{\mathcal C}(h_\partial)=\int\prod_{(v,f)}dh_{vf}\,\prod_f\delta(H_f)\prod_vA_v(\{h_{vf}\}_{f\ni v}).}
Here $H_f$ is the ordered product of wedge holonomies, including the prescribed boundary holonomy for an open face. Boundary orientations and normalizations are part of this definition.
\eq{F5}{\delta(h)=\sum_jd_j\chi_j(h).}
In an orthonormal basis on the shared graph,
\eq{F6}{W_{\mathcal C_2\circ\mathcal C_1}=\sum_{j_\Gamma,i_\Gamma}W_{\mathcal C_2}(\ldots;j_\Gamma,i_\Gamma)W_{\mathcal C_1}(j_\Gamma,i_\Gamma;\ldots).}
The common boundary is dualized according to orientation. Neither a sum over all complexes nor its continuum limit should be presented as an already established unique completion.

\section{Cartan decomposition and boosters}
\eq{B1}{g=u\,e^{r\sigma_3/2}v^{-1},\qquad u,v\in SU(2),\quad r\ge0.}
\eq{B2}{dg=\frac1{4\pi}\sinh^2r\,dr\,du\,dv.}
\eq{B3}{D^{(\rho,k)}_{jm,\ell n}(g)=\sum_{p=-\min(j,\ell)}^{\min(j,\ell)}D^j_{mp}(u)d^{(\rho,k)}_{j\ell p}(r)D^\ell_{pn}(v^{-1}).}
For normalized intertwiners, define
\eq{B4}{\boxed{\begin{aligned}
\mathcal B_4^\gamma(\mathbf j,\boldsymbol\ell;i,k)={}&\frac1{4\pi}\int_0^\infty dr\,\sinh^2r\sum_{\mathbf p}\overline{(i_i^{\mathbf j})_{\mathbf p}}(i_k^{\boldsymbol\ell})_{\mathbf p}\\
&\times\prod_{a=1}^4d^{(\gamma j_a,j_a)}_{j_a\ell_a p_a}(r).
\end{aligned}}}
If the unnormalized tensors are $q_i=i_i/\sqrt{d_i}$, then
\eq{B5}{\mathcal B_4^\gamma(\mathbf j,\boldsymbol\ell;i,k)=\sqrt{d_id_k}\,B_4^\gamma(\mathbf j,\boldsymbol\ell;i,k).}
The proposed normalized form of the gauge-fixed decomposition is
\eq{B6}{\boxed{\begin{aligned}
A_v(j_{ab},i_a)={}&\sum_{\substack{\ell_{ab}\ge j_{ab}\\2\le a<b\le5}}\sum_{k_2,\ldots,k_5}
\mathcal J_{K_5}(\ell_{ab};i_1,k_2,k_3,k_4,k_5)\\
&\times\prod_{a=2}^5\mathcal B_{4,a}^\gamma.
\end{aligned}}}
\eq{B7}{\ell_{1a}=j_{1a},\qquad a=2,3,4,5.}
There are four boosters, six independent auxiliary face-spin sums and four auxiliary intertwiners [2,5]. $\mathcal J_{K_5}$ is the central normalized tensor contraction with the prescribed orientations. It is not automatically a bare named $15j$ symbol. Before implementation, its phase and port mapping must be established from (D7)--(D9). The all-parallel booster definition (B4) must be matched to each node using the chosen duality conventions; this is part of that same check.

\section{Recoupling down to factorials}
For the conventional first-kind symbol [5],
\eq{R1}{\begin{aligned}
\begin{Bmatrix}a_1&a_2&a_3&a_4&a_5\\b_1&b_2&b_3&b_4&b_5\\c_1&c_2&c_3&c_4&c_5\end{Bmatrix}_{I}
={}&(-1)^{\sum_r(a_r+b_r+c_r)}\sum_xd_x\\
&\times\begin{Bmatrix}a_1&c_1&x\\c_2&a_2&b_1\end{Bmatrix}\begin{Bmatrix}a_2&c_2&x\\c_3&a_3&b_2\end{Bmatrix}\\
&\times\begin{Bmatrix}a_3&c_3&x\\c_4&a_4&b_3\end{Bmatrix}\begin{Bmatrix}a_4&c_4&x\\c_5&a_5&b_4\end{Bmatrix}\begin{Bmatrix}a_5&c_5&x\\a_1&c_1&b_5\end{Bmatrix}.
\end{aligned}}
This expression needs normalization factors, permutations and possibly recouplings before being identified with $\mathcal J_{K_5}$ for the supplied diagram. Equation (I8) fixes the required basis changes unambiguously.
\eq{R2}{\Delta(a,b,c)=\sqrt{\frac{(a+b-c)!(a-b+c)!(-a+b+c)!}{(a+b+c+1)!}}.}
Racah's sum is
\eq{R3}{\begin{aligned}
\begin{Bmatrix}a&b&c\\d&e&f\end{Bmatrix}={}&\Delta(a,b,c)\Delta(a,e,f)\Delta(d,b,f)\Delta(d,e,c)\\
&\times\sum_z\frac{(-1)^z(z+1)!}{(z-a-b-c)!(z-a-e-f)!(z-d-b-f)!(z-d-e-c)!}\\
&\times\frac1{(a+b+d+e-z)!(a+c+d+f-z)!(b+c+e+f-z)!}.
\end{aligned}}
Sum over exactly the integers for which all denominator factorials are nonnegative. The $3j$ endpoint is
\eq{R4}{\begin{aligned}
\begin{pmatrix}a&b&c\\m&n&p\end{pmatrix}={}&\delta_{m+n+p,0}(-1)^{a-b-p}\Delta(a,b,c)\\
&\times\sqrt{(a+m)!(a-m)!(b+n)!(b-n)!(c+p)!(c-p)!}\\
&\times\sum_z\frac{(-1)^z}{z!(a+b-c-z)!(a-m-z)!(b+n-z)!}\\
&\times\frac1{(c-b+m+z)!(c-a-n+z)!}.
\end{aligned}}

\section{Reduced Lorentz matrices and hypergeometric endpoint}
Define the phase through a norm, not a branch-ambiguous square root:
\eq{M1}{e^{i\psi_j^\rho}=\frac{\Gamma(j+1+i\rho)}{|\Gamma(j+1+i\rho)|}.}
Speziale's equation (21) [2], with $(-1)^{(j-\ell)/2}$ written as an exponential, is
\begingroup\small
\eq{M2}{\begin{aligned}
d^{(\rho,k)}_{j\ell p}(r)={}&e^{i\pi(j-\ell)/2}e^{i\psi_j^\rho-i\psi_\ell^\rho}\frac{\sqrt{d_jd_\ell}}{(j+\ell+1)!}e^{-(k-i\rho+p+1)r}\\
&\times\sqrt{(j+k)!(j-k)!(j+p)!(j-p)!(\ell+k)!(\ell-k)!(\ell+p)!(\ell-p)!}\\
&\times\sum_{s,t}\frac{(-1)^{s+t}e^{-2tr}(k+p+s+t)!(j+\ell-k-p-s-t)!}
{s!(j-k-s)!(j-p-s)!(k+p+s)!t!(\ell-k-t)!(\ell-p-t)!(k+p+t)!}\\
&\times{}_2F_1(\ell+1-i\rho,k+p+1+s+t;j+\ell+2;1-e^{-2r}).
\end{aligned}}
\endgroup
Sum over nonnegative $s,t$ for which the factorials exist. For $k=j$ one sum collapses:
\eq{M3}{\begin{aligned}
d^{(\rho,j)}_{j\ell p}(r)={}&e^{i\pi(j-\ell)/2}e^{i\psi_j^\rho-i\psi_\ell^\rho}\frac{\sqrt{d_jd_\ell}}{(j+\ell+1)!}\\
&\times\sqrt{\frac{(2j)!(\ell+j)!(\ell-j)!(\ell+p)!(\ell-p)!}{(j+p)!(j-p)!}}e^{-(j-i\rho+p+1)r}\\
&\times\sum_{s=0}^{\ell-j}\frac{(-1)^se^{-2sr}}{s!(\ell-j-s)!}\\
&\times{}_2F_1(\ell+1-i\rho,j+p+1+s;j+\ell+2;1-e^{-2r}).
\end{aligned}}
EPRL sets $\rho=\gamma j$. In the minimal block:
\eq{M4}{\boxed{d^{(\rho,j)}_{jjm}(r)=e^{-(j-i\rho+m+1)r}{}_2F_1(j+m+1,j+1-i\rho;2j+2;1-e^{-2r}).}}
\eq{M5}{{}_2F_1(a,b;c;z)=\sum_{n=0}^\infty\frac{(a)_n(b)_n}{(c)_n\,n!}z^n,\qquad(a)_n=\frac{\Gamma(a+n)}{\Gamma(a)}.}
The series is displayed on its convergence domain. Analytic continuation is a separate operation.

\section{Toller matrices in the same convention}
\eq{Q1}{\boxed{T^{(+,\rho,k)}(g)+T^{(-,\rho,k)}(g)=D^{(\rho,k)}(g).}}
\eq{Q2}{T^{(\pm,\rho,k)}_{jm,\ell n}(g)=\sum_pD^j_{mp}(u)t^{(\pm,\rho,k)}_{j\ell p}(r)D^\ell_{pn}(v^{-1}).}
To convert the analytic convention of [4] into the convention of (M2):
\eq{Q3}{D_{\rm S}=\Phi_{j\ell}(\rho)D_{\rm R},\qquad T^\pm_{\rm S}=\Phi_{j\ell}(\rho)T^\pm_{\rm R},}
\eq{Q4}{\Phi_{j\ell}(\rho)=e^{-i\pi(j-\ell)/2}\frac{\Gamma(j+1+i\rho)}{|\Gamma(j+1+i\rho)|}\frac{\Gamma(\ell+1-i\rho)}{|\Gamma(\ell+1-i\rho)|}.}
The S/R subscripts occur only in this convention card. $\Phi_{jj}=1$.
\eq{Q5}{\begin{aligned}
P_{j\ell}(\widetilde\rho;\rho)&=\prod_{n=0}^{j+\ell}\frac{i\widetilde\rho-(n-j)}{i\rho-(n-j)}\\
&=\frac{\Gamma(-j-i\rho)\Gamma(\ell+1-i\widetilde\rho)}{\Gamma(-j-i\widetilde\rho)\Gamma(\ell+1-i\rho)}.
\end{aligned}}
The Feynman representation, with the phase conversion included, is
\eq{Q6}{\boxed{\begin{aligned}
T^{(\pm,\rho,k)}_{jm,\ell n}(g)={}&\lim_{\varepsilon\downarrow0}\int_{\mathbb R}\frac{d\widetilde\rho}{2\pi i}\frac{\pm1}{\widetilde\rho-\rho\mp i\varepsilon}\\
&\times P_{j\ell}(\widetilde\rho;\rho)\frac{\Phi_{j\ell}(\rho)}{\Phi_{j\ell}(\widetilde\rho)}D^{(\widetilde\rho,k)}_{jm,\ell n}(g).
\end{aligned}}}
The phase ratio disappears for $j=\ell$. R\"uhl's analyticity statements concern $D/\Phi$ and $T/\Phi$, not the rephased functions in isolation. Potential spectral poles and distributional limits retain the prescriptions of [3,4].

For minimal blocks, with $s=\pm1$,
\eq{Q7}{\boxed{\begin{aligned}
t^{(s,\rho,j)}_{jjm}(r)={}&e^{-(j-si\rho+sm+1)r}\frac{\Gamma(2j+2)\Gamma(si\rho-sm)}{\Gamma(j-sm+1)\Gamma(j+1+si\rho)}\\
&\times{}_2F_1(j+sm+1,j+1-si\rho;1+sm-si\rho;e^{-2r}).
\end{aligned}}}
Euler transformation gives a finite polynomial. For $z=e^{-2r}$,
\eq{Q8}{\begin{aligned}
t^{(s,\rho,j)}_{jjm}(r)={}&e^{-(j-si\rho+sm+1)r}\frac{\Gamma(2j+2)\Gamma(si\rho-sm)}{\Gamma(j-sm+1)\Gamma(j+1+si\rho)}\frac1{(1-z)^{2j+1}}\\
&\times\sum_{n=0}^{j-sm}\frac{(-j-si\rho)_n(-j+sm)_n}{(1+sm-si\rho)_n\,n!}z^n.
\end{aligned}}
For $\rho\ne0$ the scalar example is
\eq{Q9}{d^{(\rho,0)}_{000}(r)=\frac{\sin(\rho r)}{\rho\sinh r},\qquad t^{(\pm,\rho,0)}_{000}(r)=\pm\frac{e^{\pm i\rho r}}{2i\rho\sinh r}.}
This is not the trivial Lorentz representation. Its separate branches cannot simply be evaluated at $\rho=0$ without a prescription.
For $j=k=m=1/2$,
\eq{Q10}{\begin{aligned}
t^+(r)&=-\frac{e^{i\rho r}}{2(\rho^2+\tfrac14)\sinh^2r},\\
t^-(r)&=\frac{e^{-i\rho r}(\cosh r+2i\rho\sinh r)}{2(\rho^2+\tfrac14)\sinh^2r}.
\end{aligned}}
Set $\rho=\gamma/2$ for EPRL. The sum is regular at $r=0$, although each term is singular.
\eq{Q11}{T^\pm(g_1g_2)\ne T^\pm(g_1)T^\pm(g_2)\qquad\text{in general}.}
This prevents incorrectly applying the Wigner factorization (D9) to a single Toller branch.

\section{Causal vertex and wedge signs}
\eq{CV1}{\sigma_a\in\{\pm1\},\qquad\kappa_{ab}=\sigma_a\sigma_b.}
Using the same contractions and orientations as (D6)--(D8),
\eq{CV2}{\begin{aligned}
A_v^{\boldsymbol\sigma}(j_{ab},i_a)={}&\int\prod_{a=2}^5dg_a\;\Contr_{K_5}\bigg[\bigotimes_{a<b}
T^{(\kappa_{ab},\gamma j_{ab},j_{ab})}_{j_{ab},j_{ab}}(g_{ab})\\
&\hspace{47mm}\otimes\bigotimes_ai_a\bigg].
\end{aligned}}
Here the displayed spin-pair subscript denotes the whole minimal magnetic block. This is a causal-model branch, not merely a more expanded EPRL formula [3].
At the integrand level,
\eq{CV3}{\prod_{a<b}D_{ab}=\sum_{\{\kappa_{ab}=\pm1\}}\prod_{a<b}T_{ab}^{\kappa_{ab}}.}
Induced signs, however, obey
\eq{CV4}{\prod_{(ab)\in C}\kappa_{ab}=1\qquad\text{for every graph cycle }C.}
On $K_5$, this gives $2^{5-1}=16$ distinct induced assignments, compared with $2^{10}=1024$ independent wedge assignments. Thus, in general,
\eq{CV5}{\sum_{\boldsymbol\sigma/\{\sigma\sim-\sigma\}}A_v^{\boldsymbol\sigma}\ne A_v^{\rm EPRL}.}
Integrated use of (CV3) requires a common regulator or sufficient convergence. Cancellation in the summed integrand does not establish finiteness of each causal sector.

For a later spectral computation branch, apply (Q6) on every wedge of a vertex with independent $\widetilde\rho_{ab}$. Interchanges of spectral integrals, vertex integrals and auxiliary-spin sums require explicit convergence assumptions; no naive $d\mapsto t$ booster replacement is proposed.

\section{Separate coherent and spinor representation mode}
\eq{CO1}{\ket{j,\mathbf n}=D^j(g_{\mathbf n})\ket{j,j},\qquad g_{\mathbf n}\hat{\mathbf z}=\mathbf n.}
\eq{CO2}{\|j_a,\mathbf n_a\rangle=\int_{SU(2)}du\,\bigotimes_aD^{j_a}(u)\ket{j_a,\mathbf n_a}.}
The semiclassical closure condition is
\eq{CO3}{\sum_aj_a\mathbf n_a=0.}
Closure characterizes geometric data; the group average in (CO2) can be defined without imposing it first.
\eq{CO4}{\ket z=\begin{pmatrix}z^0\\z^1\end{pmatrix},\quad X_i(z)=\tfrac12\bra z\sigma_i\ket z,\quad|\mathbf X(z)|=\tfrac12\langle z|z\rangle.}
\eq{CO5}{\mathbf n(z)=\frac{\bra z\boldsymbol\sigma\ket z}{\langle z|z\rangle},\qquad\ket{j,z}=\left(\frac{\ket z}{\|z\|}\right)_{\rm sym}^{\otimes2j}.}
\eq{CO6}{c_i(j_a,\mathbf n_a)=\bra i\bigotimes_a\ket{j_a,\mathbf n_a},\qquad\|j_a,\mathbf n_a\rangle=\sum_i c_i\ket i.}
\eq{CO7}{A_v(j_{ab},\mathbf n_{ab})=\sum_{\{i_a\}}A_v(j_{ab},i_a)\prod_ac_{i_a}(j_{ab},\mathbf n_{ab}).}
Dualization and coherent-state phases follow the chosen boundary orientations. A projective-spinor integral would be a further separately reviewed extension, not an unannounced implementation addition.

For appropriately phased nondegenerate Lorentzian Regge coherent boundary data, the optional asymptotic card is
\eq{CO8}{A_v(\lambda j,\mathbf n)\sim\lambda^{-12}\left[N_+e^{i\lambda s_{\rm R}}+N_-e^{-i\lambda s_{\rm R}}\right],\qquad s_{\rm R}=\gamma\sum_fj_f\Theta_f.}
$s_{\rm R}$ is the dimensionless leading Regge phase. Boundary phases and critical-point assumptions must accompany this formula. It is not a universal asymptotic statement for arbitrary spin networks.

\section{Mapping the supplied TikZ wiring to the catalogue}
The TikZ macros provide a mathematical reference for topology, port order, crossings, arrows, labels, rotation and gluing. The proposed browser format is SVG: sharp paths with semantic identifiers, permitting equation--diagram highlighting. A static TikZ rendering and downloadable source can remain available.

\begin{center}
\begin{tabular}{ll@{\qquad}ll}\toprule
Wire&Nodes&Wire&Nodes\\\midrule
E1&T4--T5&E6&T4--T2\\
E2&T3--T5&E7&T1--T4\\
E3&T2--T5&E8&T2--T3\\
E4&T5--T1&E9&T3--T1\\
E5&T4--T3&E10&T2--T1\\\bottomrule
\end{tabular}
\end{center}
The more consequential ordering is:
\begin{center}
\begin{tabular}{ll}\toprule Node&Neighbours at ports 1, 2, 3, 4\\\midrule
T1&T3, T4, T5, T2\\T2&T1, T3, T4, T5\\T3&T4, T5, T2, T1\\T4&T5, T2, T1, T3\\T5&T2, T1, T3, T4\\\bottomrule
\end{tabular}
\end{center}
Equation (I4) couples ports 1 and 2, and ports 3 and 4. Reordering them changes the recoupling convention.

The macro currently labels the internal intertwiner line with $\iota_{\cdots}$. A separate virtual-spin label $k$ is needed. Causal signs $\sigma_a$ and $\kappa_{ab}$ must also be distinguished from wire-orientation arrows. Diagram crossings are not additional graph vertices.

The remaining named-$15j$ mapping must be checked by evaluating both the magnetic-index contraction and its recoupled expression on low-spin cases. Until then, (R1) is an explicit reference convention, not a claimed direct transcription of every invocation of the macro.

\section{Source plan and review points}
The final website should attach source equation/page identifiers and every phase or normalization conversion to each formula. This draft establishes the catalogue and the source families; it does not claim a completed per-equation bibliography.

\begin{enumerate}
\item P. Martin-Dussaud, \emph{A Primer of Group Theory for Loop Quantum Gravity and Spin-foams}, supplied PDF. \href{https://arxiv.org/abs/1902.08439}{arXiv:1902.08439}. Basis, generators, duality, graphical calculus and representation theory.
\item S. Speziale, \emph{Boosting Wigner's nj-symbols}, supplied PDF. \href{https://arxiv.org/abs/1609.01632}{arXiv:1609.01632}. In particular equations (11)--(13), (21)--(25), and the convention appendices.
\item E. Bianchi, C. Chen and M. Gamonal, \emph{Causal spinfoam vertex for 4d Lorentzian quantum gravity}, supplied PDF. \href{https://arxiv.org/abs/2601.23162}{arXiv:2601.23162}. Causal signs, vertex and Feynman kernel.
\item E. Bianchi, C. Chen and M. Gamonal, \emph{Toller matrices and the Feynman $i\varepsilon$ in spinfoams}, supplied PDF. \href{https://arxiv.org/abs/2604.24945}{arXiv:2604.24945}. Phase conversion, analytic convention, minimal Toller blocks and spectral transform.
\item P. Don\`a and P. Frisoni, \emph{How-To compute EPRL spin foam amplitudes}. \href{https://arxiv.org/abs/2202.04360}{arXiv:2202.04360}. Gauge fixing, booster factorization and the first-kind $15j$ expansion.
\item C. Rovelli and F. Vidotto, \emph{Covariant Loop Quantum Gravity}. Classical/canonical background, compact presentation, quantum tetrahedra and coherent states. \href{https://www.cpt.univ-mrs.fr/~rovelli/IntroductionLQG.pdf}{Author-hosted manuscript}.
\end{enumerate}

The low-spin matrices (L7) and (L9) were checked directly in the Clebsch--Gordan basis (I4). Equations (G6), (L8)--(L11) and the terminating-polynomial transformation (Q8) are derived identities. The most useful initial review targets are the common phase convention, normalized versus unnormalized intertwiners, endpoint orientations, the tetrahedral volume normalization, and the scope of the causal branch.
\end{document}
